% ==============================================================================
% S1 Introduction
% ==============================================================================

\section{Introduction}
\label{sec:introduction}

This paper examines how a coordinate-dependent scoring framework
can produce reproducible phase-localized alignment patterns in a
geographically clustered corpus. The corpus is UNESCO Cultural and Mixed
World Heritage sites ($N = \NclusterTotal$).
A global parameter sweep over frequency and phase selects a
high-scoring frame near 3\textdegree{} spacing and the
$\sim$35\textdegree{}E corridor (\S\ref{sec:results_unit_sweep},
\S\ref{sec:sensitivity_global}). All subsequent analyses are conditional
on this selected frame. Because the frame is chosen by optimizing over the
data, results at this frame reflect post-selection and should not be read
as independent statistical evidence. Gerizim provides one historically
documented point within the identified corridor
(\S\ref{sec:anchor}).
Within this frame, the data show weak but reproducible departures from
spatial uniformity. The same pattern is absent at non-optimal anchors and
is not recovered as a spectral peak under invariant periodogram analysis.

Three diagnostic checks point in the same direction: dome
von~Mises mixture, Spearman cluster asymmetry, and phase-rotation null
(Table~\ref{tab:bonferroni_both}; $p_\mathrm{adj} = \pHolmTestOne$,
$\pHolmTestTwo$, $\pHolmTestThree$). Effects are small (A+ rate
\clusterApRate\% vs.\ \NullRateAp{} null; $\rho = \spearmanDevClusterRho$);
the anomaly localizes the $\sim$35\textdegree{}E corridor, not a unique
anchor. This is compatible with 3\textdegree{} spacing but does not
establish it
(\S\ref{sec:anchor}, \S\ref{sec:results_unit_sweep}).

The observed spatial regularity is concentrated in the longitude range
where Eurasian cultural site density is highest (the Eurasian
high-density corridor, often referred to as the Silk Road region) and
is consistent with the geographic clustering of culturally related sites.
The Wikidata Q180987 stupa inventory and OWTRAD route network show
similar enrichment patterns under the same reference frame, but share
geographic and cultural structure with UNESCO. Their agreement is
cross-corpus consistency, not independent confirmation
(\S\ref{sec:transmission}).

%------------------------------------------------------------------------------
\section{Background}
\label{sec:prior}
%------------------------------------------------------------------------------



\subsection{Prior Work on Spatial Periodicity and Sacred Landscapes}

Statistical evaluation of spatial periodicity in monument distributions has a
long and contested history \citep{thom1967,kendall1974,freeman1976}.
Probability-based frameworks for archaeoastronomical alignments have been
advanced by Magli \citep{magli2013,magli2016} and Ruggles and Silva
\citep{ruggles1999astronomy,ruggles2015handbook,silva2020statistical};
the ICOMOS-IAU thematic study \citep{ruggles2010icomos} provides institutional
context for World Heritage analysis.

\subsection{Reference Phase and Candidate Anchors}
\label{sec:anchor}

The global anchor sweep (\S\ref{sec:sensitivity_global}) places the
$\sim$35\textdegree{}E corridor among the highest-scoring reference phases
under the harmonic-deviation score, with marginal discovery-corrected support:
both Gerizim and Jerusalem place the
35.25\textdegree{}E modulo-3\textdegree{} phase band in the global top \anchorSweepTopPct\% of
\anchorSweepNanchors{} trial anchors, while Mecca ($p \approx \pMeccaAnchor$,
\pMeccaAnchorSig{}), Mt.\ Meru ($p \approx \pMeruAnchor$, \pMeruAnchorSig{}),
and Mt.\ Kailash ($p \approx \pKailashAnchor$, \pKailashAnchorSig{}) yield
no enrichment. Harmonic deviations are computed from Mount Gerizim
(\GerizimLon\textdegree{}E), which lies within that corridor.
The Samaritan \textit{tabbur ha-aretz} (``navel of the earth'') tradition
\citep{purvis1968,macdonald1964,crown1989} provides independent historical
motivation for this choice. Tests using Gerizim reflect properties of the
selected corridor, not of any uniquely identified meridian.

Gerizim is not inscribed (UNESCO tentative ref.~5706, DMS \GerizimLonDMS{});
Jerusalem is inscribed and self-excluded when used as anchor.


%------------------------------------------------------------------------------
\section{Data and Methods}
\label{sec:methods}
%------------------------------------------------------------------------------

\subsection{Corpus Construction}
\label{sec:corpus}

The UNESCO World Heritage XML dataset \citep{unescowhc2024xml} contains
$\NwhcTotal$ inscribed nominations. UNESCO designates each site
\textit{Cultural} (criteria i to vi), \textit{Natural} (criteria vii to x),
or \textit{Mixed}; sites with cultural heritage alongside natural values
receive Mixed and are included here.
Retaining Cultural and Mixed ($-\NexcNatural$ Natural-only) and
excluding $\NexcNoCoord$ entries lacking geocoordinates (both linear
route networks) yields $N = \NclusterTotal$ sites.

\subsection{Geodetic Framework}
\label{sec:geodetic}

The reference anchor is Mount Gerizim (\GerizimLon\textdegree{}E; not
inscribed, \S\ref{sec:anchor}). Jerusalem is inscribed and scored as a
regular corpus site. For sensitivity sweeps, Gerizim is added synthetically
so each anchor self-excludes symmetrically.

For a site at longitude $\lambda$, the deviation from the nearest 3\textdegree{} harmonic is
$\delta = |\lambda - \GerizimLon| \bmod 3^\circ$ (folded to $[0, 1.5^\circ]$).
The tier classification is a fixed angular criterion.
Kilometer labels (e.g.\ $\lesssim\ApThreshKm{}$~km for Tier-A+) are equatorial
approximations ($\delta \times 111$~km/\textdegree{}). No latitude correction is applied.
Kilometer-scale distances should be interpreted as approximate: UNESCO coordinates
are site centroids, and precision below $\sim$1~km is not meaningful.

The geometric null for Tier-A+ is $2 \times \ApThreshDeg{}^\circ / 1.5^\circ = \NullRateApPct{}\%$
(assuming uniform longitudes). A Gaussian KDE gives a similar expected A+
count under longitude clumping ($\simKDEMean \pm \simKDEStd$ expected A+).
The primary diagnostics use simulation-based nulls (\S\ref{sec:simulation_null}).


\subsection{Statistical Framework}
\label{sec:statistical_framework}

The choice of $(\omega, \phi) = (3^\circ, \text{Gerizim})$ follows
from a global parameter sweep.
Define $T = \max_{\omega,\phi} S(\omega, \phi)$ where $S$ is the
harmonic-deviation enrichment score; $T$ is calibrated by permuting site
longitudes and recomputing, correcting for the full search.
Global scans over $\phi$ (\anchorSweepNanchors{} anchors) and unit scans
over $\omega$ both maximize at $3^\circ$ / $35^\circ$E
(\S\ref{sec:results_unit_sweep}, \S\ref{sec:sensitivity_global};
anchor-sweep permutation $p = \anchorMaxPermP$~\anchorMaxPermPSig{}).
$S$ tests lattice proximity under a specified frame, not intrinsic
periodicity; the two are epistemically distinct (\S\ref{sec:repn_dependence}).
Three threshold-free diagnostics below characterize alignment at this fixed frame
(see \S\ref{sec:test_populations}):

\begin{description}
  \item[\textbf{Test~1}] Two-component von~Mises/uniform mixture fit to the
    dome sub-corpus ($N = \NcircTotal$), mean fixed at Gerizim phase, assessed
    by parametric bootstrap.
  \item[\textbf{Test~2}] Spearman rank correlation between harmonic node mean
    deviation and site count, full corpus ($N = \NclusterTotal$), permutation null.
  \item[\textbf{Test~3}] Phase-rotation null: the reference grid is shifted by a
    uniform random offset in $[0^\circ,3^\circ)$; tests whether the Gerizim phase
    minimizes the full-corpus mean harmonic deviation.
\end{description}

All three point in the same direction (diagnostic $p$-values in
Table~\ref{tab:bonferroni_both}; Holm-adjusted at $k = \BonfK$).
All primary diagnostics operate on the same harmonic-deviation representation
and should be interpreted as complementary summaries of a single
constructed feature, not as independent evidence.
These $p$-values quantify structure within the constructed frame and
are not evidence against spatial randomness in an invariant sense.
Tier labels (A++, A+, A) are descriptive vocabulary only and do not enter
any diagnostic statistic (\S\ref{sec:tier_classification}).

\subsection{Tier Classification}
\label{sec:tier_classification}

Sites are assigned to tiers by their angular deviation $\delta$ from the
nearest 3\textdegree{} harmonic node.
The three nested A-windows span the same range of proximities examined
in the multi-scale enrichment analysis (\S\ref{sec:multiscale}):
4\% null ($\approx$\AppThreshKm{}~km) through 20\% null ($\approx$\ATierThreshKm{}~km),
covering the full interval over which the four corpora show enrichment.
These windows are not identified as peaks; they are fixed proximity bands
used as descriptive vocabulary for summarising harmonic proximity.
No primary test statistic depends on a tier cutoff.

\begin{description}
  \item[\textbf{Tier~A++}]
    $\delta \leq \AppThreshDeg{}^\circ$ ($\lesssim\AppThreshKm{}$~km); $P_0 = \NullRateApp{}$.
  \item[\textbf{Tier~A+}]
    $\delta \leq \ApThreshDeg{}^\circ$ ($\lesssim\ApThreshKm{}$~km); $P_0 = \NullRateAp{}$.
  \item[\textbf{Tier~A}]
    $\delta \leq \ATierThreshDeg{}^\circ$ ($\lesssim\ATierThreshKm{}$~km); $P_0 = \NullRateA{}$.
  \item[\textbf{Tier~B}] Middle zone, between harmonics and inter-harmonic midpoints.
  \item[\textbf{Tier~C, C$-$, C$--$}] Symmetric mirror of A, A+, and A++, measured from the
    nearest inter-harmonic midpoint.
\end{description}


\subsection{Test Populations}
\label{sec:test_populations}

\subsubsection*{Primary threshold-free diagnostics (Tests~1 to~3)}

Tests~1 to~3 are described in full in \S\ref{sec:statistical_framework}.

The only additional detail specific to this section is the keyword
population for Test~1: \NcircKeywords{} morphological keywords from four
roots (\textit{stupa/stupas}, \textit{tholos}, \textit{dome/domed/domes},
\textit{spherical}) are matched as whole words against each site's name,
XML description, and extended OUV text ($N = \NcircTotal$).
The context-validated subset ($N = \NcircValidated$) restricts to sites
where the keyword appears in an architectural context and is reported as
a robustness check (Test~5x).

\subsubsection*{Supporting enrichment tests (Tests~4 to~6)}
\label{sec:enrichment_tests}

Tests~4 to~6 count sites falling within a fixed proximity window of the
nearest harmonic node and test whether the rate exceeds a geometric null.
Three nested windows are reported (Tiers~A++, A+, A; null rates \NullRateApp{}, \NullRateAp{}, \NullRateA{};
widths \AppThreshDeg{}$^\circ$, \ApThreshDeg{}$^\circ$, \ATierThreshDeg{}$^\circ$ on $[0^\circ, 1.5^\circ)$).
No primary diagnostic depends on a tier cutoff (see \S\ref{sec:tier_classification}).


\textbf{Test~4 (full-corpus enrichment).}
Full corpus ($N = \NclusterTotal$), no keyword filter.
Binomial test of A+ rate vs.\ $p_0 = \NullRateAp{}$ null.

\textbf{Test~5 (dome enrichment).}
Dome keyword subset ($N = \NcircTotal$).
Fisher exact and binomial tests of A+ and A++ rates.

\textbf{Test~6 (mound evolution, sensitivity).}
Extends the dome subset with \textit{tumulus, tumuli, barrow, barrows, kofun}
($N = \NevoTotal$); not in the primary diagnostic family.

\subsubsection*{Descriptive and post-hoc tests}

\textbf{Test~7 (temporal gradient).}
Cochran-Armitage trend test on inscription year. Descriptive only.

\textbf{Test~8 (founding classifier, post-hoc).}
Keyword classifier; no weight in the primary evidential chain
(Appendix~\ref{sec:kwdaudit_founding}).

\subsection{Multiple-Comparisons Protocol}
\label{sec:bonferroni}

\textbf{Primary diagnostic family ($k = \BonfK$, Holm step-down).}
Tests~1 to~3 are all threshold-free; they form the primary diagnostic family.
The adjustment calibrates these summaries within the selected representation.
It is not a confirmatory familywise test of discovery.

\textbf{Supporting enrichment tests (Tests~4 and~5).}
Tests~4 and~5 use a pre-specified proximity window and are not in the primary diagnostic family.
Results are in the FDR audit (Table~\ref{tab:bonferroni_both}; \S\ref{sec:fdraudit}).

\textbf{Other.}
Test~6 is a sensitivity variant of Test~5.
Test~7 (temporal gradient) is descriptive.
Test~8 is post-hoc.
Religion sub-group analyses (\S\ref{sec:resultsrelig}) and diagnostic
Tests~A to~D (\S\ref{sec:results_rayleigh}) are not Bonferroni-corrected.

\subsection{Simulation Null Models}
\label{sec:methods_simulation}

Three simulation approaches (\simNperms{} trials each), a within-dome
bootstrap stability check, and a restricted geographic-concentration null
are described with their results in \S\ref{sec:simulation_null}.

\subsection{Computational Reproducibility}
\label{sec:scripts}

All statistics are generated by an automated pipeline; every script writes
to a shared results store and LaTeX macros are regenerated without manual
editing. Scripts, audit files, and reproduction instructions are catalogued
in Supplementary~S1; see also \S\ref{sec:data_code}.

%------------------------------------------------------------------------------
\section{Results}

\label{sec:results}

Effects are weak in absolute terms (full-corpus A+ rate \clusterApRate\%
against a \NullRateAp{} null; Spearman $\rho = \spearmanDevClusterRho$).

%------------------------------------------------------------------------------
\subsection{Spacing Specificity: Is 3\textdegree{} Discriminated from Nearby Intervals?}
\label{sec:results_unit_sweep}
%------------------------------------------------------------------------------

The 3\textdegree{} spacing is the highest-scoring parameter from the
unit sweep (\S\ref{sec:statistical_framework}); metrological context is
noted in \S\ref{sec:anchor}. Within the canonical range examined, no nearby random spacing
reproduces the observed tier-bin peak: a longitude-preserving permutation
of the dome sub-corpus (\permSlopeSpecN{} trials) yields
$p = \permSlopeCanonBestJointDomeAP{}$ at A$^{+}$ (\permSlopeCanonBestJointSigDomeAP{})
and $p = \permSlopeCanonBestJointDomeAPP{}$ at A$^{++}$
(\permSlopeCanonBestJointSigDomeAPP{}). Restricting to morphologically
unambiguous spherical forms (stupa, tholos, $n = \permSlopeKwN{}$) gives
$p = \permSlopeCanonBestJointKwA{}$ (\permSlopeCanonBestJointSigKwA{}) at
the A tier.

A $\pm 0.3\%$ shift from exactly 3\textdegree{} collapses this peak in at
least one population, so the result is best read as compatibility with a
3\textdegree{} model rather than precise recovery of a 3\textdegree{}
unit (\S\ref{sec:metrology}).

The 1.5\textdegree{} and 6\textdegree{} spacings produce similar tier-bin
counts but are mathematical consequences of the 3\textdegree{} structure
(shared harmonic nodes). Neither constitutes an independent signal.
Candidate intervals at 4\textdegree{} and 8\textdegree{} (\S\ref{sec:limitations_geo})
are similarly arithmetic consequences of the 3\textdegree{} grid.
Results support relative optimality of 3\textdegree{} within the tested
family. It is a representative optimal scale under this scoring function,
not a uniquely identified spacing. Other spacings produce comparable
localized structure under alternative parameterizations.


%------------------------------------------------------------------------------
\subsection{Phase localization}
\label{sec:results_phase_localization}
%------------------------------------------------------------------------------

The 3\textdegree{} model performs better under the tested diagnostics than
adjacent alternatives (\S\ref{sec:results_unit_sweep}); we next ask whether the
pattern is uniformly distributed over all possible phases or concentrated at a
specific phase.
A phase-envelope scan over all origins $\phi \in [0,3^\circ)$
(\phaseNphases{} phases at \phaseStep\textdegree{} resolution) shows a strong
preferred phase: full-corpus A+ counts span
\phaseFullMin{} to \phaseFullMax{} (median \phaseFullMedian{}), and Gerizim's phase
($\phaseGerizimPhase^\circ$ mod 3\textdegree{}) sits at the
\phaseFullGerPctile\textsuperscript{th} percentile (\phaseFullGerCount{} A+),
the \phaseDomeGerPctile\textsuperscript{th} for domes
(\phaseDomeGerCount/\NcircTotal{}; range \phaseDomeMin{} to \phaseDomeMax{}).
Rayleigh permutation on the 3\textdegree{}-circle:
$p_\mathrm{perm} = \rayleighPermP$~\rayleighPermPSig{}
(\S\ref{sec:phase_peak_formal}).
This phase was identified as part of the selected frame (post-selection;
\S\ref{sec:statistical_framework}).

\subsection{Cluster Asymmetry (Tests~2 and~3, Primary)}
\label{sec:resultscluster}

Tests~2 and~3 use the full corpus ($N = \NclusterTotal$), no keyword filter
(\S\ref{sec:test_populations}), and quantify alignment without any tier cutoff
(Table~\ref{tab:bonferroni_both}).
\textbf{Test~2} (Spearman rank correlation between harmonic node mean deviation and site
count): $\rho = \spearmanDevClusterRho$, permutation $p = \spearmanDevClusterP$
(\spearmanDevClusterPSig{}). Nodes drawing more sites show smaller mean harmonic
deviations.
\textbf{Test~3} (phase-rotation null): the observed full-corpus mean harmonic deviation
is smaller than expected when the reference grid is rotated by a uniform random offset
in $[0^\circ, 3^\circ)$ ($p_\mathrm{perm} = \fullMeanDevPermP$, \fullMeanDevPermPSig{}).
Both point in the same direction (Holm-adjusted diagnostics:
$p_\mathrm{adj} = \pHolmTestTwo$~\pHolmTestTwoSig{};
$p_\mathrm{adj} = \pHolmTestThree$~\pHolmTestThreeSig{}).

The tier-based characterisation below provides descriptive context.
For each site, count other UNESCO sites within $\pm \clusterRadiusDeg$\textdegree{}
longitude; compare cluster sizes for A+ vs.\ non-A+ sites (Mann-Whitney $U$,
permutation test). At the harmonic level, compare per-harmonic site counts
for A+-bearing vs.\ non-A+ harmonics.

\begin{table*}[!t]
\centering
\caption{Cluster asymmetry: A+ vs.\ non-A+ sites in the full UNESCO corpus
  ($N = \NclusterTotal$; cluster radius $\pm \clusterRadiusDeg$\textdegree{} lon)}
\label{tab:clusterasym}
\small
\adjustbox{max width=\textwidth}{%
\begin{tabular}{p{6.2cm}rrrr}
\toprule
Metric & A+ sites & Non-A+ sites & Ratio / $Z$ & $p$ \\
\midrule
Mean cluster size (sites within $\pm \clusterRadiusDeg$\textdegree{} lon)
  & \clusterApMean & \clusterNonApMean & \clusterRatio$\times$ & \\
Mann-Whitney $U$ (one-tailed)
  & & & & \clusterMWp~\clusterMWpSig{} \\
Permutation (\NpermCluster{} shuffles)
  & & & $Z = \clusterPermZ$ & \clusterPermP~\clusterPermPSig{} \\
\midrule
\multicolumn{5}{l}{\textit{Harmonic-level density (all Tier $\leq$ B sites per 3\textdegree{} harmonic):}} \\
Harmonics with $\geq 1$ A+ site (\NclusterApHarmonics{})
  & \multicolumn{2}{c}{mean $= \clusterHarmonicApMean$ sites}
  & & \\
Harmonics with 0 A+ sites (\NclusterNonApHarmonics{})
  & \multicolumn{2}{c}{mean $= \clusterHarmonicNonApMean$ sites}
  & \clusterHarmonicRatio$\times$ & $\clusterHarmonicMWpStr$~\clusterHarmonicMWpSig{} \\
\midrule
\bottomrule
\end{tabular}}% end adjustbox
\end{table*}

\noindent A+-bearing harmonics host \clusterHarmonicRatio$\times$ more sites
on average than non-A+ harmonics ($p\ \clusterHarmonicMWpStr$;
Table~\ref{tab:clusterasym}).
Sites at harmonic nodes cluster more densely than sites at harmonic midpoints
(\clusterNodeApMean{} vs.\ \clusterAntiApMean{} neighbors;
$\clusterNodeVsAntiApRatio\times$, $p = \clusterNodeVsAntiApMWp$).

\paragraph*{Subtest: does clustering produce the unit-sensitivity peak?}
\label{sec:cluster_cond_sensitivity}
A joint permutation test (\clusterCondN{} iterations) finds the conditional
fraction of clustering-matched permutations producing the canonical-unit peak
(\clusterCondPeakFrac{}) is identical to the unconditional rate
(\clusterCondPeakFracUncond{}): in this check, the clustering and
unit-sensitivity summaries appear separable.

\paragraph*{Methodological note: Rayleigh phase-lock (Test~B).}
\label{sec:results_rayleigh}
The near-perfect Rayleigh $R$ of the A+ sub-sample is a mathematical
consequence of the A+ classification criterion (\S\ref{sec:circularity})
and is reported as a structural description.
Three formal permutation tests (\NpermRayleigh{} shuffles each, seed~\simSeed{})
address distinct aspects of the
periodicity:

\textbf{Test~A (full-corpus Rayleigh, negative control).}
Full-corpus Rayleigh $R = \fullRayleighR$
($Z = \fullRayleighZ$, $p_\mathrm{perm} = \fullRayleighPermP$, \fullRayleighPermPSig{}):
no spurious global 3\textdegree{} periodicity; enrichment is concentrated in A+ sites.

\textbf{Test~B (A+-only Rayleigh, structural check).}
Among \Nap{} Tier-A+ sites, $R = \rayleighR$ ($Z = \rayleighZ$):
a structural consequence of the A+ criterion, not independent evidence
(\S\ref{sec:circularity}).

\textbf{Tests~C to~E (anchor specificity).}
Circular-shift null: \anchorShiftApCount{} A+ sites vs.\ null mean \anchorShiftNullMu{} ($p_\mathrm{perm} = \anchorShiftPermP$, \anchorShiftPermPSig{}).
Targeted von~Mises (full corpus, $\mu_0$ fixed at Gerizim phase): $\hat{w} = \targetedWfull$ (${\approx}\targetedWfullPct\%$ of sites), $p_\mathrm{boot} = \targetedPfull$ (\targetedPfullSig{}).
Global anchor sweep ($\anchorSweepNanchors{}$ trial anchors): Levantine corridor $p = \anchorMaxPermP$ (\anchorMaxPermPSig{}).

\subsection{Full-Corpus Enrichment (Test~4, Supporting)}
\label{sec:results_global}

We now quantify the excess at the Gerizim-grid phase.

\begin{figure*}[!t]
\centering
\includegraphics[width=0.9\textwidth]{../figures/fig_devhist}
\caption{Distribution of harmonic deviations for the full UNESCO corpus.
  The excess of sites within the Tier-A+ window ($\delta \leq \ApThreshDeg{}^\circ$,
  shaded) above the uniform null (dashed line) illustrates the corpus-level
  enrichment pattern; Tests~1 and~2 (domed monuments and cluster asymmetry)
  characterize its morphological and spatial structure.}
\label{fig:devhist}
\end{figure*}

The full UNESCO Cultural/Mixed corpus ($N = \NclusterTotal$) shows a
Tier-A+ rate of \clusterApRate\% (Figure~\ref{fig:devhist})
(\NclusterAp/\NclusterTotal; Clopper-Pearson 95\% CI
[\clusterApCIlo\%, \clusterApCIhi\%]),
above the \NullRateAp{} null
(binomial $p = \clusterApBinom$; supporting Test~4, not in the primary
diagnostic family).
The observed enrichment is $\clusterEnrichAp\times$, modest relative
to the dome subpopulation ($\circEnrichAp\times$) and the harmonic cluster
asymmetry. The primary threshold-free diagnostics are reported in
Table~\ref{tab:bonferroni_both}; they summarize structure within the
selected harmonic representation, not independent confirmation.

\textbf{Region-stratified audit.}
A region-specific empirical null (Fisher combination of per-region binomials)
yields $p = \pRegionEmpAll$ (\pRegionEmpAllSig{}), so it does not provide
region-stratified evidence against geographic concentration. The pre-2000
cohort remains elevated against the geometric null
($p = \pPreTwoK$~\pPreTwoKSig{}; \S\ref{sec:resultstemporal}), and is treated
as descriptive context.

\subsubsection*{World religion sites}
\label{sec:resultsrelig}
Religion sub-populations are defined by keyword sets documented in the supplementary audit archive~\citep{tenenbaum2026audit}.
A chi-squared test finds no significant A+ rate difference across the five main
religion groups ($\chi^2 = \religChiSqStat$, df\,=\,\religChiSqDof,
$p = \pReligChiSq$, \pReligChiSqSig{}).
Both Buddhism (\NbudATier/\NbudSitesRelig{} A-tier, \budATierRate\%;
\NbudInterHarmonic/\NbudSitesRelig{} C-band, \budInterHarmonicRate\%;
$p = \pBudJoint$~\pBudJointSig{}) and Hinduism (\NhinATier/\NhinSites{} A-tier, \hinATierRate\%;
\NhinInterHarmonic/\NhinSites{} C-band, \hinInterHarmonicRate\%;
$p = \pHinJoint$~\pHinJointSig{}) show a joint A-tier/C-band signal
(multivariate hypergeometric). Both persist in tradition-exclusive sites
($p = \pBudExclJoint$~\pBudExclJointSig{}, $p = \pHinExclJoint$~\pHinExclJointSig{}).

%------------------------------------------------------------------------------
\subsection{Domed and Spherical Monuments (Test~5, Supporting)}
\label{sec:resultscircular}

The keyword filter is described in \S\ref{sec:test_populations} ($N = \NcircTotal$).
The dome sub-corpus should be treated as a hypothesis-generating subpopulation:
the keyword set was refined iteratively during analysis, and the enrichment
should not be read as confirmatory evidence independent of that process.
\textbf{Test~1} (dome von~Mises mixture, Table~\ref{tab:bonferroni_both}): a two-component
von~Mises/uniform mixture with the mean fixed at the Gerizim phase yields
$\hat{w} = \targetedWdome$ ($\hat{\kappa} = \targetedKdome$),
$p_\mathrm{boot} = \targetedPdome$ (\targetedPdomeSig{};
Holm $p_\mathrm{adj} = \pHolmTestOne$, \pHolmTestOneSig{}).
Approximately $\targetedWdomePct\%$ of dome sites are concentrated at the Gerizim phase,
with an angular spread of ${\approx}\targetedSdDomeDeg$\textdegree{}.
The tier-based enrichment below provides complementary descriptive context.

\textbf{Tier-A++ concentration (supporting result).}
\NcircTierApp{} of \NcircTotal{} dome sites (\NcircAppRate\%, $\circEnrichApp\times$ the \NullRateApp{} null)
fall within Tier-A++ ($\delta \leq \AppThreshDeg{}^\circ$, $\lesssim\AppThreshKm{}$~km).
This concentration is directionally reproducible across three test statistics:
Fisher OR~$= \circAppFisherOR$, $p = \pCircAppFisher$~\pCircAppFisherSig{};
permutation $Z = \circAppPermZ$, $p = \circAppPermP$~\pCircAppPermSig{};
bootstrap $p = \circAppBootP$~\pCircAppBootSig{}.
Dome A++ gives consistent results across all three null models and remains directionally robust under full Bonferroni at $k = \NtotalTests$ (\S\ref{sec:multiple_comparisons}).

\textbf{Tier-A+ enrichment (supporting Test~5).}
\NcircTierAp/\NcircTotal{} sites at Tier-A+
(\NcircApRate\% [\circApCIlo, \circApCIhi], $\circEnrichAp\times$;
Fisher OR~$= \circApFisherOR$, $p = \pCircApFisher$~\pCircApFisherSig{};
binomial $p = \pCircAp$~\pCircApSig{}).
Tier-A: \NcircTierA/\NcircTotal{} (\circARate\%, $\circEnrichA\times$;
Fisher $p = \pCircAFisher$~\pCircAFisherSig{};
binomial $p = \pCircA$~\pCircASig{}).
The $\chi^2$-uniform test ($p = \pCircChi$, \pCircChiSig{}) is consistent with
concentration at harmonic nodes.
Tier-B (middle zone): \NcircTierB/\NcircTotal{} sites;
mean $\delta = \circMeanDev$\textdegree{}.

\noindent\textbf{Tier-A+ sites}:
Khoja Ahmed Yasawi (0.2~km),
Boyana Church (0.3~km),
Lumbini (0.8~km),
Humayun's Tomb (2.0~km),
Pyu Ancient Cities (2.3~km),
Silk Roads: Chang'an-Tianshan (2.5~km), Trulli of Alberobello (3.6~km),
Old City of Jerusalem (4.1~km),
Borobudur (4.3~km),
Kizhi Pogost (5.0~km), Echmiatsin/Zvartnots (6.2~km).

A leave-one-out binomial check shows that no single site drives the result:
removing any A+ site yields $N = \LOOdomeN$, A+$= \LOOdomeAp$,
$p = \LOOdomeP$.
Leave-2-out (worst case) across all $\binom{\LKOnAp}{2}$ combinations:
$p = \LKOworstTwoP$.
Leave-3-out (worst case): $p = \LKOworstThreeP$; the evolution corpus (Test~6) (\S\ref{sec:resultsevolution_inline}) 
and the A++ concentration are more robust. 
The A++ leave-3-out worst case across all
$\binom{\LKOAppNApp}{3} = \LKOAppNCombosThreeApp$ triples yields
$p = \LKOAppworstThreeP$, remaining below $\alpha = 0.05$ in this audit
(leave-2-out worst case $p = \LKOAppworstTwoP$).

\textbf{Multi-scale combined enrichment.}
\label{sec:multiscale}
A Fisher combination across \multiscaleK{} log-spaced proximity windows
(null rates 0.33\% to 40\%; canonical tiers A++, A+, A are three of the \multiscaleK{})
shows all windows enriched above their null (OR 1.2$\times$ to 6.7$\times$,
14 of 15 individually $p < 0.10$);
permutation-calibrated $p = \multiscaleDomeCombinedP$~\multiscaleDomeCombinedPSig{}.
No fixed tier cutoff enters this statistic.
Applied jointly across all four corpora (UNESCO domes, UNESCO full,
Wikidata Q180987 stupas, OWTRAD Silk Road vertices), the maximum joint enrichment
falls at approximately \multiscaleJointMaxKm{} km from a harmonic node
(1/10th of the \textasciitilde{}289 km inter-node spacing at 30\textdegree{}N),
with permutation-calibrated $p = \multiscaleJointMaxP$~\multiscaleJointMaxPSig{}.
Each corpus peaks within this \multiscaleJointMaxKm{} km bound
(4 to 29 km across corpora), a scale that emerges from the data rather than any
pre-specified threshold.

\textbf{Context-validated robustness (Test~5x).}
Sentence-level validation ($N = \NcircValidated$, \NcircContextRejected{} rejected):
A++ $p = \pCircAppValidatedFisher$~\pCircAppValidatedFisherSig{};
A+ $p = \pCircApValidatedFisher$~\pCircApValidatedFisherSig{}.
The A++ result strengthens when false positives are removed; A+ remains
directionally elevated under this stricter filter.
Keyword-stratified check: stupa/stupas alone ($n = \kwSensStupaN$) carries the strongest signal
(A++ $p = \kwSensStupaAppP$~\kwSensStupaAppPSig{});
dome/tholos/spherical without stupa ($n = \kwSensNoStupaN$) remains elevated
($p = \kwSensNoStupaAppP$~\kwSensNoStupaAppPSig{}).

\textbf{Geographic-concentration checks (dedicated simulation).}
\label{sec:geo_conc_null}
Domed monuments are concentrated in the Eurasian corridor:
\domeEurasianFraction\% fall in 20\textdegree{} to 120\textdegree{}E
versus \fullCorpusEurasianFraction\% of the full corpus.
Two targeted resampling checks evaluate whether this explains the dome enrichment
(\simNperms{} trials).

\textbf{Null~A (within-dome bootstrap stability check):} Draw $N = \NcircTotal$ longitudes
from the dome sites' own list. Mean $= \geoNullDomeBootMean$ A+,
$p = \geoNullDomeBootP$ (\geoNullDomeBootPSig{}). The dome sites' positions are concentrated
near harmonic longitudes; this check characterizes the observed dome
distribution rather than serving as an external geographic null.

\textbf{Null~C (restricted geographic-concentration null):} Pool all full-corpus sites
within $\pm W$\textdegree{} of any dome site and bootstrap the expected A+
count. Dome sites (\NcircTierAp~A+, \NcircApRate\%) exceed
the broader footprint at all three windows tested:
$\pm 2$\textdegree{} ($p = \geoNullDomeRestrictedPtwo$, \geoNullDomeRestrictedPtwoSig{});
$\pm 5$\textdegree{} ($p = \geoNullDomeRestrictedP$, \geoNullDomeRestrictedPSig{});
$\pm 10$\textdegree{} ($p = \geoNullDomeRestrictedPten$, \geoNullDomeRestrictedPtenSig{}).
Results are stable across the tested window widths.


\begin{figure*}[!t]
\centering
\includegraphics[width=0.9\textwidth]{../figures/fig_null_c}
\caption{Null~C bootstrap distributions for the dome sub-population
  ($N = \NcircTotal$, observed A+ $= \NcircTierAp$) at three
  window widths. For each window $\pm W$\textdegree{}, all full-corpus
  sites within $\pm W$\textdegree{} of any dome longitude form the
  restricted pool; $10^5$ draws of size $N = \NcircTotal$ from that
  pool define the expected A+ distribution (blue bars). The observed
  count (red line) exceeds the pool mean by $Z \approx 2.7$ in all three
  windows ($\pm 2$\textdegree{}: $p = \geoNullDomeRestrictedPtwo$, \geoNullDomeRestrictedPtwoSig{};
  $\pm 5$\textdegree{}: $p = \geoNullDomeRestrictedP$, \geoNullDomeRestrictedPSig{};
  $\pm 10$\textdegree{}: $p = \geoNullDomeRestrictedPten$, \geoNullDomeRestrictedPtenSig{}),
  showing stability across the tested window widths.}
\label{fig:null_c}
\end{figure*}

In the global anchor sweep (\S\ref{sec:sensitivity_global}),
Gerizim scores \GerizimSweepAp{} A+ (\anchorSweepPctile{}\textsuperscript{th}
percentile); Jerusalem scores \JerusalemSweepAp{}
(\JerusalemSweepPctile{}\textsuperscript{th}).

\subsection{Hemispherical Mound Evolution (Test~6, Sensitivity)}
\label{sec:resultsevolution_inline}

Extending the dome keyword subset of Test~5 with earthen mound-stage keywords (see \S\ref{sec:test_populations});
the combined population ($N = \NevoTotal$) adds
\NevoMoundOnly{} mound-stage sites to the core
domed/spherical population, with \NevoOverlap{}
sites matching both keyword sets.

\noindent\textbf{Stage-by-stage breakdown} (Table~\ref{tab:evostages1b}):

\begin{table*}[!t]
\centering
\caption{Test~6 enrichment by evolutionary stage ($N = \NevoTotal$).
  Null rates: A++~$P_0 = \NullRateApp{}$, A+~$P_0 = \NullRateAp{}$, A~$P_0 = \NullRateA{}$.
  Fisher exact tests: stage vs.\ full corpus background.}
\label{tab:evostages1b}
\small
\adjustbox{max width=\textwidth}{%
\begin{tabular}{lrrllrllrl}
\toprule
 & & \multicolumn{2}{c}{Tier~A++ ($\leq\AppThreshKm{}$\,km)} & & \multicolumn{2}{c}{Tier~A+ ($\leq\ApThreshKm{}$\,km)} & & \multicolumn{2}{c}{Tier~A ($\leq\ATierThreshKm{}$\,km)} \\
\cmidrule(lr){3-4}\cmidrule(lr){6-7}\cmidrule(lr){9-10}
Stage & $N$ & A++ & Rate OR ($p$) & & A+ & Rate OR ($p$) & & A & Rate OR ($p$) \\
\midrule
Mound\textsuperscript{a} & \NevoMound
  & \NevoMoundApp & \evoMoundAppRate\% $\evoMoundAppFisherOR\times$ (\pEvoMoundAppFisher~\pEvoMoundAppFisherSig{}) &
  & \NevoMoundAp  & \evoMoundApRate\%  $\evoMoundFisherOR\times$ (\pEvoMoundFisher~\pEvoMoundFisherSig{}) &
  & \NevoMoundA   & \evoMoundARate\%   $\evoMoundAFisherOR\times$ (\pEvoMoundAFisher~\pEvoMoundAFisherSig{}) \\
Stupa\textsuperscript{b} & \NevoStupa
  & \NevoStupaApp & \evoStupaAppRate\% $\evoStupaAppFisherOR\times$ (\pEvoStupaAppFisher~\pEvoStupaAppFisherSig{}) &
  & \NevoStupaAp  & \evoStupaApRate\%  $\evoStupaFisherOR\times$ (\pEvoStupaFisher~\pEvoStupaFisherSig{}) &
  & \NevoStupaA   & \evoStupaARate\%   $\evoStupaAFisherOR\times$ (\pEvoStupaAFisher~\pEvoStupaAFisherSig{}) \\
Dome\textsuperscript{c}  & \NevoDome
  & \NevoDomeApp  & \evoDomeAppRate\%  $\evoDomeAppFisherOR\times$ (\pEvoDomeAppFisher~\pEvoDomeAppFisherSig{}) &
  & \NevoDomeAp   & \evoDomeApRate\%   $\evoDomeFisherOR\times$ (\pEvoDomeFisher~\pEvoDomeFisherSig{}) &
  & \NevoDomeA    & \evoDomeARate\%    $\evoDomeAFisherOR\times$ (\pEvoDomeAFisher~\pEvoDomeAFisherSig{}) \\
\midrule
\textbf{Combined} & \textbf{\NevoTotal}
  & \textbf{\NevoApp} & \textbf{\evoAppRate\%} $\mathbf{\evoAppFisherOR\times}$ (\textbf{\pEvoAppFisher~\pEvoAppFisherSig{}}) &
  & \textbf{\NevoAp}  & \textbf{\evoApRate\%}  $\mathbf{\evoApFisherOR\times}$ (\textbf{\pEvoApFisher~\pEvoApFisherSig{}}) &
  & \textbf{\NevoA}   & \textbf{\evoARate\%}   $\mathbf{\evoAFisherOR\times}$ (\textbf{\pEvoAFisher~\pEvoAFisherSig{}}) \\
\bottomrule
\multicolumn{10}{p{1.2\linewidth}}{\footnotesize
  \textsuperscript{a}tumulus/barrow/kofun.\;
  \textsuperscript{b}stupa/stupas.\;
  \textsuperscript{c}dome/tholos/spherical.
  OR = odds ratio; $p$ = one-sided Fisher exact vs.\ full UNESCO cultural corpus ($N = \NclusterTotal$).
  Stage $N$ values sum to more than the combined total because \NevoOverlap{} sites match multiple stages.} \\
\end{tabular}}
\end{table*}

\noindent The combined population ($N = \NevoTotal$, A+ = \evoApRate\%;
Table~\ref{tab:evostages1b}) is consistent with Test~1. The mound stage does not reach
significance (\evoMoundApRate\%, OR $= \evoMoundFisherOR\times$, Fisher $p = \pEvoMoundFisher$, \pEvoMoundFisherSig{}). The signal is carried
by the stupa (A+ \evoStupaApRate\%, OR $= \evoStupaFisherOR\times$, $p = \pEvoStupaFisher$~\pEvoStupaFisherSig{}) and dome
sub-stages (A+ \evoDomeApRate\%, OR $= \evoDomeFisherOR\times$, $p = \pEvoDomeFisher$~\pEvoDomeFisherSig{}).
Leave-one-out checks confirm stability.
Leave-3-out (worst case, $\EvoLKOnCombosThree{}$ triples):
$p = \EvoLKOworstThreeP$.

\textbf{Stupa geographic-concentration null.}
\label{sec:stupa_geo_null}
Stupa sites are concentrated in South/Southeast Asia
(\stupaHeartlandFrac\% in 70\textdegree{} to 110\textdegree{}E).
Null~A (within-stupa bootstrap): $p = \stupaGeoBootP$ (\stupaGeoBootPSig{}).
Null~C (restricted $\pm 5$\textdegree{} draw): $p = \stupaGeoRestrictedP$ (\stupaGeoRestrictedPSig{}).

\subsection{Temporal Gradient (Test~7, Descriptive)}
\label{sec:resultstemporal}

The Tier-A++ rate in the founding canon (1978 to 1984) is \cohortAppArate\%
(\NcohortAppAapp/\NcohortAppA{}; $\canonEnrich\times$ the \NullRateApp{} null,
$p = \pCanonApp$~\pCanonAppSig{}), declining monotonically toward the null
in later cohorts (2000 to 2025: \cohortAppErate\%).
Cochran-Armitage three-cohort trend test (A++, decreasing):
$Z = \ZcochranAppThree$, $p = \pCochranAppThree$~(\pCochranAppThreeSig{}).
The five-cohort trend is marginal ($p = \pCochranAppFive$~\pCochranAppFiveSig{}).
By contrast, the A+ rate shows no significant decreasing trend
($p = \pCochranThree$~\pCochranThreeSig{}), consistent with the gradient
being specific to the highest-precision tier.
This is consistent with UNESCO's founding canon over-representing
historically prominent monuments, some of which are also precisely aligned
to harmonic nodes.
The Mantel-Haenszel stratified test is underpowered at the sub-group level;
Logistic regression with region as a covariate is directionally consistent
with the A++ gradient
($Z = \zLogRegCohortApp$, $p = \pLogRegCohortApp$~\pLogRegCohortAppSig{},
OR~$= \orLogRegCohortApp$).
The declining rate may partly reflect UNESCO's expansion into regions where hemispherical traditions are less prevalent.
The gradient is treated as descriptive context only (Figure~\ref{fig:temporal}).

\begin{figure*}[!t]
\centering
\includegraphics[width=0.9\textwidth]{../figures/fig_temporal}
\caption{Tier-A++ and Tier-A+ rates by UNESCO inscription-year cohort.
  The founding canon (1978 to 1984) shows an elevated A++ rate
  (\cohortAppArate\%, $p = \pCanonApp$~\pCanonAppSig{}) that declines
  monotonically across later cohorts toward the \NullRateApp{} null (dotted line;
  Cochran-Armitage $Z = \ZcochranAppThree$, $p = \pCochranAppThree$~\pCochranAppThreeSig{}).
  The A+ rate (right bars) shows no significant trend ($p = \pCochranThree$~\pCochranThreeSig{}).}
\label{fig:temporal}
\end{figure*}

\begin{table*}[tp]
\centering\small
\caption{Multiple-comparisons summary.
  Primary diagnostic family: Tests~1 to~3 (threshold-free, $k = \BonfK$).
  Supporting enrichment tests (Tests~4 and~5) use a pre-specified $p_0 = \NullRateAp{}$
  proximity window; see \S\ref{sec:enrichment_tests}.}
\label{tab:bonferroni_both}
\adjustbox{max width=\textwidth}{%
\begin{tabular}{lcrll}
\toprule
Test & $N$ & Raw $p$ & Holm $p_\mathrm{adj}$, $k{=}\BonfK$ & Status \\
\midrule
Test~1 (dome VM mix.)  & \NcircTotal     & \targetedPdome$^{b}$ ($w = \targetedWdome$, $\kappa = \targetedKdome$) & \pHolmTestOne~\pHolmTestOneSig{}  & \textbf{Diagnostic} \\
Test~2 (Spearman)      & \NclusterTotal  & \spearmanDevClusterP$^{b}$ ($\rho = \spearmanDevClusterRho$) & \pHolmTestTwo~\pHolmTestTwoSig{}  & \textbf{Diagnostic} \\
Test~3 (mean dev.)     & \NclusterTotal  & \fullMeanDevPermP$^{b}$ & \pHolmTestThree~\pHolmTestThreeSig{} & \textbf{Diagnostic} \\
\midrule
Test~4 (full-corpus A+) & \NclusterTotal & \clusterApBinom$^{a}$ & N/A  & Supporting \\
Test~5 (dome A+)        & \NcircTotal    & \pCircAp$^{a}$ (OR~$= \circApFisherOR$) & N/A & Supporting \\
Test~6 (mound evol.)   & \NevoTotal      & \pEvoApFisher$^{a}$ (OR~$= \evoApFisherOR$) & N/A  & Sensitivity \\
\midrule
Test~7 (temporal)      & \NoriginCorpus  & \pCochranThree$^{b}$ & N/A  & Descriptive \\
\bottomrule
\multicolumn{5}{p{0.95\textwidth}}{\footnotesize
  $^{a}$Binomial or Fisher exact $p$ vs.\ $\NullRateAp{}$ geometric null (window $= \ApThreshDeg{}^\circ$; convenience choice, sensitivity-audited over $2\%$ to $10\%$ null rates in Supplementary~S1);
  $^{b}$permutation or bootstrap $p$.
  Tests~1 to~3 are threshold-free; no tier cutoff enters any primary diagnostic.
  Adjusted $p$-values describe calibration within the selected representation,
  not confirmatory rejection of a global discovery null.
  Holm (1979) step-down \citep{holm1979simple}.} \\
\end{tabular}}
\end{table*}

\subsection{FDR audit and Bonferroni correction}
\label{sec:fdraudit}

Results are consolidated in Table~\ref{tab:bonferroni_both}.
A comprehensive FDR audit (\NtotalTests{} tests;
\citealt{benjamini1995controlling}) retains \NsurviveFDR{} at $q < 0.05$;
\NsurviveBonfAll{} remain below the full-Bonferroni audit threshold at $k = \NtotalTests$:
dome/spherical A++ ($p = \BonfAllSurvivepCircApp$),
cluster permutation ($p = \BonfAllSurviveclusterPermP$),
harmonic Mann-Whitney ($p = \BonfAllSurviveclusterHarmonicMWp$),
sequential-peak permutation ($p = \BonfAllSurvivepeakSigP$),
and evolution combined A++ ($p = \BonfAllSurvivepEvoApp$).

\section{Robustness Checks}
\label{sec:robustness}

\subsection{Unit Sweep and Sensitivity Slope}
\label{sec:metrology}

Unit sweep and sensitivity slope results are in \S\ref{sec:results_unit_sweep}.
Briefly: nearby non-harmonic spacings are non-significant; a $\pm 0.3\%$
shift collapses joint significance; Monte Carlo sweep
(\S\ref{sec:limitations_geo}) supports the relative peak;
Tests~C and~D (\S\ref{sec:results_rayleigh}) support phase specificity
within the selected 3\textdegree{} frame.

\subsection{Global Anchor Sweep}
\label{sec:sensitivity_global}

The global sweep (\anchorSweepNanchors{} trial anchors, 0.01\textdegree{}
resolution) uses the full corpus plus synthetic Gerizim
(\S\ref{sec:geodetic}), with each anchor self-excluding
($N = \NgerizimSweepCorpus$ per trial).

Both anchors place the corridor in the global top \anchorSweepTopPct\%
(\anchorSweepNanchors{} trial anchors, 0 to 360\textdegree{}E;
Gerizim: \GerizimSweepAp{} A+, \anchorSweepPctile\textsuperscript{th} percentile;
Jerusalem: \JerusalemSweepAp{} A+, \JerusalemSweepPctile\textsuperscript{th}
percentile). The result is anchor-invariant within the
\GerizimJerusalemKm\,km corridor window.
Other historical meridians tested here do not show comparable enrichment.

\subsubsection*{The 35.25\textdegree{}E modulo-3\textdegree{} phase band}
\label{sec:phase_peak_artifact}

Every longitude congruent to 35.25\textdegree{}E modulo 3\textdegree{}
(\NharmonicsCoarse{} anchors)
produces identically the maximum A+ count, spaced 3\textdegree{} apart.
The flat-top width is a function of the tier threshold chosen:
A++ ($\pm\AppThreshDeg{}^\circ$), A+ ($\pm\ApThreshDeg{}^\circ$), A ($\pm\ATierThreshDeg{}^\circ$);
the number of tied anchors is a methodological artifact, not a physical feature.
Mount Gerizim is a fixed landmark at \GerizimLon\textdegree{}E (phase $\GerizimPhase$\textdegree{},
gap $\phaseGap$\textdegree{} $\approx \phaseGapKm$~km). It lies within the A and A+ flat-top
bands but not A++.

\subsubsection*{Formal statistical tests for the 35.25\textdegree{}E modulo-3\textdegree{} phase band}
\label{sec:phase_peak_formal}

The full-corpus Rayleigh ($p_\mathrm{perm} = \fullRayleighPermP$, \fullRayleighPermPSig{}) shows no
global 3\textdegree{} spectral periodicity; Test~B is tautological.
Tests~C ($p_\mathrm{perm} = \anchorShiftPermP$, \anchorShiftPermPSig{}) and~D ($p_\mathrm{boot} = \targetedPfull$, \targetedPfullSig{})
are positive anchor-specific tests; Test~E ($p = \anchorMaxPermP$, \anchorMaxPermPSig{}, marginal) is the anchor discovery correction.
The unit sensitivity sweep (\S\ref{sec:metrology}) and Monte Carlo unit sweep
(\S\ref{sec:limitations_geo}) identify 3\textdegree{} as the highest-scoring
spacing under the template score, \emph{given} the pre-specified Gerizim anchor;
neither sweep constitutes anchor-free or period-free discovery.

\textbf{Corridor and landmark ranking.}
Of the \corridorUniqueGer{} A+ sites unique to Gerizim and
\corridorUniqueJer{} unique to Jerusalem, \corridorSharedAp{} are shared;
the enrichment belongs to the corridor, not to either anchor individually.
Using each UNESCO site as anchor, Gerizim ranks
\#\GerizimTempleRank{} among the extended candidate landmarks
(\GerizimPctileAp\textsuperscript{th} percentile).

\subsection{Spatial Autocorrelation}
\label{sec:autocorrelation}

A+ harmonics host $\clusterHarmonicRatio\times$ as many total UNESCO
sites as non-A+ harmonics ($p\ \clusterHarmonicMWpStr$).
DEFF correction via ICC ($\pm 0.5$\textdegree{}, $\rho = \ICChalf$):
DEFF~$= \DEFFhalf$, $N_\mathrm{eff} = \NeffHalf$,
corrected $p = \pNeffHalf$.
Block bootstrap (3\textdegree{} blocks, \simNperms{} resamples):
95\% CI $[\blockBootCIlo\%, \blockBootCIhi\%]$ excludes the \NullRateAp{} null.
Within-group permutation is non-discriminating for alignment vs.\ geographic
concentration; the restricted geographic-concentration null
(\S\ref{sec:geo_conc_null}) addresses that comparison directly.

\subsection{Null Model Validation}
\label{sec:simulation_null}

Three simulation nulls (\simNperms{} trials each, seed~$= \simSeed{}$,
\citealt{good2005permutation}):
(1)~Longitude permutation: dome ($N = \NcircValidated$)
$p_\mathrm{perm} = \simDomePermP$~\simDomePermPSig{};
1978-1984 canon ($N = \NcanonSites$) $p = \simCanonPermP$~\simCanonPermPSig{};
pre-2000 ($N = \NpreTwoK$) $p = \simPreTwoKPermP$~\simPreTwoKPermPSig{}.
(2)~Bootstrap from full corpus: dome $p_\mathrm{boot} = \simDomeBootP$~\simDomeBootPSig{}.
(3)~KDE-smoothed null: expected $\simKDEMean \pm \simKDEStd$, observed \Nap{}, $Z = \simKDEZ$,
$p = \simKDEP$~\simKDEPSig{}.
The KDE null ($p = \simKDEP$~\simKDEPSig{}) supports dome enrichment; permutation and bootstrap
are marginal at the A+ level but stronger for the A++ sub-population
(permutation $p = \circAppPermP$~\pCircAppPermSig{},
bootstrap $p = \circAppBootP$~\pCircAppBootSig{}).
The within-dome bootstrap characterizes the observed concentration near
harmonic longitudes, while the restricted geographic-concentration null
(Null~C; \S\ref{sec:geo_conc_null}) indicates that longitude footprint
alone is not sufficient to explain the dome enrichment.

\subsection{Pipeline Portability: Wikidata Q180987 Stupa Corpus}
\label{sec:stupas_q180987}

The pipeline was applied to $N = \wikiStupaTotal{}$ Wikidata-attested stupas
(class Q180987) in a separately assembled corpus \citep{tenenbaum2026audit}.
The Kedu Plain (Java) hosts the densest cluster at the 75\textdegree{} harmonic (75\textdegree{}E from Gerizim),
anchored by Borobudur (c.\ 800~CE, \wikiStupaBorobudurDevKm~km, Tier~A+)
\citep{miksic1990borobudur}; the World Peace Pagoda at Lumbini achieves
A++ precision ($\sim$\wikiStupaLumbiniDevM~m) under the Gerizim reference
frame.

Using the identical anchor and null, the corpus yields
\wikiStupaATierCount/\wikiStupaTotal{} Tier-A sites
(\wikiStupaATierRate\%, $\wikiStupaATierEnrich\times$ null;
binomial $p\ \wikiStupaATierBinomP$, \wikiStupaATierBinomStar;
permutation $Z = \wikiStupaPermZ$, $p_\mathrm{perm} = \wikiStupaPermP$,
\wikiStupaPermStar).
Tier-A+ is nominally elevated under the simple binomial count
(\wikiStupaApCount/\wikiStupaTotal{}, \wikiStupaApRate\%,
binomial $p = \wikiStupaApBinomP$~\wikiStupaApBinomStar), but the
portable signal is strongest at Tier-A because this external-corpus
analysis ranks cluster-level occupancy rather than individual-site
precision. Its tier emphasis should therefore not be read as mechanically
identical to the UNESCO dome scoring.
Harmonic nodes with at least one Tier-A site carry $\wikiStupaClusterRatio\times$
more stupa sites on average than empty nodes, mirroring the cluster-asymmetry
result of Test~3.

\textbf{Geographic tier breakdown and bimodal structure.}
All regional p-values below are one-sided Fisher exact tests.
A sub-regional breakdown reveals a bimodal pattern consistent with the
Buddhist and Hindu bimodal enrichment in the UNESCO religion audit (\S\ref{sec:resultsrelig}).
The Java/Sumatra node (105\textdegree{} to 115\textdegree{}E, near the 75\textdegree{} harmonic at 110.3\textdegree{}E)
concentrates \wikiJavaNodeATierN{} of \wikiJavaNodeN{} sites at Tier-A (\wikiJavaNodeATierRate\%, OR $= \wikiJavaNodeATierOR\times$, $p~\wikiJavaNodeATierP$~\wikiJavaNodeATierPSig{}),
rising to \wikiJavaTightATierN/\wikiJavaTightN{} (\wikiJavaTightATierRate\%, OR $= \wikiJavaTightATierOR\times$, $p~\wikiJavaTightATierP$~\wikiJavaTightATierPSig{}) in the 107\textdegree{} to 112\textdegree{}E window,
where A+ is also elevated at \wikiJavaTightApRate\% (OR $= \wikiJavaTightApOR\times$, $p\ \wikiJavaTightApP$~\wikiJavaTightApPSig{}).
The physical heartland, India/Nepal (70\textdegree{} to 90\textdegree{}E) and Myanmar/Cambodia
(90\textdegree{} to 105\textdegree{}E), is A-tier depleted (combined OR $= \wikiHeartlandATierOR\times$,
$p~\wikiHeartlandATierP$~\wikiHeartlandATierPSig{}, one-sided depletion test)
but strongly enriched in the inter-harmonic C-band:
Myanmar/Cambodia reaches \wikiMyanmarCbandRate\% (OR $= \wikiMyanmarCbandOR\times$, $p~\wikiMyanmarCbandP$~\wikiMyanmarCbandPSig{}) and the
combined heartland \wikiHeartlandCbandRate\% (OR $= \wikiHeartlandCbandOR\times$, $p~\wikiHeartlandCbandP$~\wikiHeartlandCbandPSig{}).
This A/C anti-correlation mirrors the bimodal pattern in the religion audit (\S\ref{sec:resultsrelig}).
Non-heartland Wikidata stupas (outside 70\textdegree{} to 110\textdegree{}E,
$N = \wikiNonHeartlandN{}$) show the same A-enriched, C-depleted profile:
\wikiNonHeartlandATierRate\% reach Tier-A (OR $= \wikiNonHeartlandATierOR\times$, $p~\wikiNonHeartlandATierP$~\wikiNonHeartlandATierPSig{})
with no C-band enrichment (OR~$= \wikiNonHeartlandCbandOR\times$, $p = \wikiNonHeartlandCbandP$~\wikiNonHeartlandCbandPSig{}).

\subsection{Multiple Comparisons}
\label{sec:multiple_comparisons}

See \S\ref{sec:bonferroni} for the full protocol.
Gerizim ranks at the \anchorSweepPctile{}\textsuperscript{th} percentile;
a $\pm 0.3\%$ unit shift collapses joint significance.
FDR audit (\NtotalTests{} tests) retains \NsurviveFDR{} at $q < 0.05$;
\NsurviveBonfAll{} remain below the Bonferroni audit threshold at $k = \NtotalTests$.

\section{Limitations}
\label{sec:limitations}

\subsection{Exploratory Status and Multiple-Comparisons Exposure}
\label{sec:limitations_circularity}
\label{sec:circularity}

The analysis is exploratory: spacing, phase, anchor, sub-corpus, and tier
thresholds were all examined iteratively on a public dataset for which
prospective constraint is not applicable. The Holm correction covers only
the three primary diagnostics; $p$-values are descriptive, not confirmatory.
They are not tests of discovery. They quantify how sharply the selected
representation describes the data.
Key degrees of freedom (3\textdegree{} spacing, Gerizim phase, keyword set,
tier thresholds) have external motivation; threshold sensitivity is in
Supplementary~S1. External corpora are not independent of UNESCO
(\S\ref{sec:transmission}).

\subsection{Geographic and Cultural Selection Bias}
\label{sec:limitations_geo}

The UNESCO list over-represents Europe and the Mediterranean-to-South-Asian
corridor (Europe/N.~America \CorpusPctEurope\%; Arab States \CorpusPctArabStates\%;
Asia-Pacific \CorpusPctAsiaPac\%; Latin America \CorpusPctLatAm\%; Africa \CorpusPctAfrica\%).
Domed monuments exceed the A+ rate of other monuments from the same longitude
footprint (Null~C, $p = \geoNullDomeRestrictedP$, \geoNullDomeRestrictedPSig{}),
which indicates that architectural type may contribute to the enrichment beyond
longitude footprint alone.

A Monte Carlo unit sweep (\mcNperms{} permutations, \mcNcandidates{} candidate
spacings 0.5\textdegree{} to 15\textdegree{}) gives the highest template score
at 3\textdegree{} ($p < 0.001$), with harmonics at 6\textdegree{} and
12\textdegree{} also elevated under the same score.
The dome subset yields \mcDomeNsigCollapsed{} independent high-scoring spacing
($P = \mcDomeCollapsedPEqOne{}$ under the geographic null).

\subsection{Representation Dependence}
\label{sec:repn_dependence}

The observed structure depends on the representation. It is recovered under
nearest-harmonic-deviation scoring but not by invariant spectral methods
(no Rayleigh peak at 3\textdegree{} in any sub-corpus).
Three converging diagnostics: (i)~a $\pm 0.3\%$ unit perturbation
collapses joint significance (\S\ref{sec:results_unit_sweep});
(ii)~multi-scale enrichment is broad, not a sharp spectral line;
(iii)~the anchor-sweep maximum is marginal ($p = \anchorMaxPermP$,
\anchorMaxPermPSig{}) while targeted tests at the identified frame
are stronger (\S\ref{sec:sensitivity_global}).
The pattern appears under the representation, not as intrinsic
periodicity. It arises from the interaction between the scoring function
and the imposed coordinate frame.
A clustered control shows that the unconstrained global maximum over
all $(\omega,\phi)$ is typical under realistic dome-site clustering
($p = \clusterCtrlPermP$, $N_{\mathrm{sim}} = \clusterCtrlNsim$) and does
not uniquely recover the hypothesised frame.
The free-sweep maximizer returns $\omega = \clusterCtrlObsOmega$\textdegree{},
a spacing with no metrological or waypoint interpretation, confirming
that the unconstrained maximum reflects spatial clustering rather than
any meaningful periodicity.
This outcome is expected: the enrichment statistic uses a fixed absolute
threshold (0.15\textdegree{}), so large spacings receive an artificially
low null rate and therefore inflate the enrichment ratio regardless of
any real structure. The free sweep is uninformative precisely because it
is not scale-invariant; the metrologically motivated prior $\omega = 3$\textdegree{}
is tested conditionally, not recovered by unconstrained search.

\subsection{Americas Negative Control}
\label{sec:americas_negative_control}
The UNESCO Americas sub-corpus ($N = \AmericasN$) shows A+ rate
\AmericasApRate\%, below the \NullRateAp{} null
($p = \AmericasOneSidedP$~\AmericasOneSidedStar{}, one-sided depletion).

\subsection{Other Limitations}

\textbf{Geocoding precision.}
UNESCO XML coordinates are administrative centroids (unsystematic noise, no directional bias).

\textbf{Metrological evidence.}
The 3\textdegree{} template was initially motivated as a base-10 subdivision
of the 30\textdegree{} b\={e}ru \citep{hunger2019mul,powell1987metrological}.
No ancient surveying instrument calibrated to such sub-units has been
recovered, and the metrological parallel is treated as heuristic motivation,
not evidence. All thresholds are degree-denominated.

%------------------------------------------------------------------------------
\section{Discussion}
\label{sec:discussion}

\subsection{Template Alignment vs.\ Intrinsic Periodicity}
\label{sec:template_vs_intrinsic}

The present results are best treated as \emph{template alignment}
(\S\ref{sec:repn_dependence}): a reproducible, phase-localized pattern
appears under the representation but is not invariant and does not imply a
periodic generative model. The analysis does not demonstrate an underlying
spatial law or intentional placement. It shows that coordinate-aligned
statistics can produce reproducible deviations in clustered datasets.
Interpretation should focus on properties of the representation, not the
existence of a latent periodic structure.

\subsection{Alternative Explanations}
\label{sec:alt_explanations}

\textbf{(1) Independent convergence.}
Hemispherical structures recur across cultures
\citep{aveni2001skywatchers,aveni2003archaeoastronomy}, but convergence
predicts no longitude specificity.

\textbf{(2) Geographic concentration.}
Null~C: other monuments from the same longitude footprint are less
enriched than the dome subset ($p = \geoNullDomeRestrictedP$,
\geoNullDomeRestrictedPSig{}); these checks reduce, but do not eliminate,
the possibility that enrichment partly reflects corridor concentration.

\textbf{(3) Random coincidence.}
Simulation nulls (KDE, permutation, bootstrap) support dome enrichment within
the selected representation,
most strongly at A++ (permutation $p = \circAppPermP$~\pCircAppPermSig{};
full results in \S\ref{sec:simulation_null}).
Tests~C and~D support phase specificity within the selected frame
(\S\ref{sec:results_rayleigh}).

\subsection{Cross-Corpus Consistency: OWTRAD and Wikidata}
\label{sec:transmission}

Two additional corpora (the Wikidata Q180987 stupa inventory and the
OWTRAD route network \citep{ciolek2004owtrad}) are analyzed under the
same harmonic-deviation representation. These datasets are not statistically
independent of the UNESCO corpus, as they share geographic and cultural
structure. Their role here is not replication, but to assess whether the
observed alignment behavior is stable across related datasets under a
common representation.

\textbf{OWTRAD route network.} The OWTRAD route network contains
\NOwtradEdges{} route segments and \NOwtradVertices{} unique endpoints.
Under the same representation, the degree-weighted endpoint distribution
shows A++ enrichment (\NOwtradDegWApp{} of \NOwtradDegW{} endpoint positions,
$z = \zOwtradDegWApp$, $p\ \pOwtradDegWApp$~\pOwtradDegWAppSig{}), while the deduplicated endpoint
distribution is weaker (\NOwtradVertexApp{}/\NOwtradVertices{}, $p =
\pOwtradVertexApp$~\pOwtradVertexAppSig{}). A phase-shift test gives
$p = \pOwtradVertexAppPhase$~\pOwtradVertexAppPhaseSig{}, and a
cluster-asymmetry check gives $p = \pOwtradClustDeg$~\pOwtradClustDegSig{}.
These results are treated as shared-representation consistency, not
independent replication. The Silk Road corridor is illustrated in
Figure~\ref{fig:geo_trail}; harmonic tier distributions for both OWTRAD
populations are shown in Figure~\ref{fig:owtrad_tiers}.

\begin{figure*}[!t]
  \centering
  \includegraphics[width=\textwidth]{../figures/fig_owtrad_tiers}
  \caption{OWTRAD Silk Road endpoint tier distributions:
  \NOwtradVertices{} deduplicated endpoints (left) and \NOwtradDegW{} degree-weighted
  route endpoint positions (right).
  The degree-weighted distribution shows A++ enrichment
  ($z = \zOwtradDegWApp$, $p\ \pOwtradDegWApp$~\pOwtradDegWAppSig{});
  the deduplicated endpoint distribution is weaker
  ($z = \zOwtradVertexApp$, $p = \pOwtradVertexApp$~\pOwtradVertexAppSig{}).}
  \label{fig:owtrad_tiers}
\end{figure*}

\begin{figure*}[!t]
  \centering
  \includegraphics[width=\textwidth]{../figures/fig_geo_trail}
  \caption{Tier distribution of UNESCO domed/stupa and Wikidata Q180987
  stupa sites along the Eurasian corridor.
  Overland route edges from the OWTRAD attested-route database
  \citep{ciolek2004owtrad}.
  Maritime route is a schematic coastal polyline after
  \citet{casson1989periplus}, \citet{sen2014silk}, and \citet{tomber2008rome}.
  An interactive version is provided as supplementary material
  (\texttt{supplementary/interactive\_map/silkroad\_interactive.html}).}
  \label{fig:geo_trail}
\end{figure*}

\section{Conclusion}

This paper reports an exploratory pattern: the UNESCO
Cultural and Mixed corpus shows a consistent template-alignment excess near
the $\sim$35\textdegree{}E corridor under multiple diagnostic checks
(Table~\ref{tab:bonferroni_both}; $p_\mathrm{adj} = \pHolmTestOne$,
$\pHolmTestTwo$, $\pHolmTestThree$), with small effect sizes and
corridor-level rather than point localization. These $p$-values describe
within-frame consistency, not confirmatory evidence.

A multi-scale combined test reduces dependence on any
single tier choice (dome subset $p = \multiscaleDomeCombinedP$~\multiscaleDomeCombinedPSig{};
joint enrichment peaks near \multiscaleJointMaxKm{}~km,
$p = \multiscaleJointMaxP$~\multiscaleJointMaxPSig{}).
The Wikidata stupa corpus and OWTRAD route network show similar enrichment
patterns under the same reference frame (\S\ref{sec:transmission}).

In short, a harmonic-alignment statistic identifies a weak but reproducible
pattern under a specific coordinate representation.
This pattern is not invariant, does not support a periodic generative model,
and is consistent with the interaction between geographic clustering and
the imposed harmonic frame. Its primary contribution is methodological. It
shows how representation choice can induce stable but non-invariant spatial
patterns in culturally clustered datasets.
Independent re-examination on corpora without UNESCO's Eurasian
sampling structure is the appropriate next step.


\appendix

\section{Von Mises Mixture Characterisation of External Corpora}
\label{app:vonmises}

A von~Mises + uniform mixture model
$f(\theta|\kappa) = w \cdot \mathrm{VM}(0,\kappa) + (1-w)\cdot\mathrm{Uniform}(0,2\pi)$
was fitted to the harmonic phase deviations ($\theta \in [0,\pi]$, $\theta=0$
on-harmonic) of two external corpora using maximum likelihood with a
grid of starting values.

OWTRAD Silk Road nodes ($N = \vmNOwtrad{}$):
$\hat{\kappa} = \vmKappaOwtrad{}$, $\hat{w} = \vmMixtureWeightOwtrad{}$,
$\hat{\sigma} = \kappa^{-1/2} = \vmSigmaDegOwtrad{}^\circ$.

Wikidata Q180987 stupas ($N = \vmNStupa{}$):
$\hat{\kappa} = \vmKappaStupa{}$, $\hat{w} = \vmMixtureWeightStupa{}$,
$\hat{\sigma} = \vmSigmaDegStupa{}^\circ$.

$N$-weighted pooled $\hat{\sigma} = \vmSigmaDegPooled{}^\circ$;
$\kappa$ ratio = \vmKappaRatio{} across corpora.

These values characterize external-corpus clustering and are not used to
derive tier boundaries; agreement with tier widths is descriptive context.

\section{Tier-Threshold Sensitivity Audit}
\label{app:tier_sensitivity}

A log-scale null-rate sweep across \logSweepNtaus{} thresholds supports
the interpretation that enrichment is not an artifact of the specific tier width:
\logSweepNsig{} of \logSweepNtaus{} thresholds reach $p < 0.05$ for the
dome subset; dome OR ranges \logSweepMinEnrichCanon{}$\times$--\logSweepMaxEnrichCanon{}$\times$
across the three canonical tiers (A, A+, A++).
Within $\pm$1 log-decade of the A+ threshold,
\logSweepAregionNsig{} of \logSweepAregionNtaus{} thresholds are significant.
Full results in \texttt{supplementary/audit/tier\_sensitivity\_audit.txt}.


%------------------------------------------------------------------------------
\section{Keyword Audit Summary}
\label{sec:kwdaudit}

Full site-by-site audit files are indexed in Supplementary~S1
(\path{supplementary/Supplementary_S1.md}; \citealt{tenenbaum2026audit});
regenerate all with \texttt{bash} \path{tools/update_all_audits.sh}.

\subsection{Dome/Spherical Monument Filter (Tests~1, 5, 5x, 6x)}
\label{sec:kwdaudit_dome}

\NcircKeywords{} morphological keywords from four roots
(\textit{stupa/stupas}, \textit{tholos}, \textit{dome/domed/domes},
\textit{spherical}) are matched as whole words against each site's
name, XML description, and extended OUV text. All \NcircTotal{} matching
sites are included in the Test~1 dome mixture and Test~5 tier-enrichment
analyses without context gating.

For Test~5x, ambiguous keywords (\textit{dome, domed, domes, spherical})
are subject to a two-gate context check: Gate~1 (\NdomeGateNeg{} negative
patterns) rejects non-architectural uses; Gate~2 (\NdomeGatePos{} positive
patterns) requires architectural confirmation. Unambiguous keywords
(\textit{stupa/stupas}, \textit{tholos}) are accepted without validation.
The seven
rejected sites are: Hiroshima Peace Memorial, Richtersveld, Gyeongju,
Vi\~{n}ales Valley, Rock Islands Southern Lagoon, Uluru-Kata Tju\d{t}a,
and Diqu\'{i}s Stone Spheres.

\hyperref[sec:resultsevolution_inline]{Test~6} adds earthen mound-stage
keywords (\textit{tumulus, tumuli, barrow, barrows, kofun} unambiguous;
\textit{mound} context-validated) to trace the hemispherical tradition
from burial mound through stupa to architectural dome. Test~6x applies
the same architectural context gates to the extended population.
Statistics for all four variants are reported in
\hyperref[sec:resultscircular]{\S\ref*{sec:resultscircular}} and
\hyperref[sec:resultsevolution_inline]{\S\ref*{sec:resultsevolution_inline}}.

\subsection{Founding-Origin Classifier (Post-Hoc)}
\label{sec:kwdaudit_founding}

\textbf{Results.} The classifier finds \foundKwApRate\% of A+ sites
carry a founding keyword versus \foundKwBaseRate\% in the full corpus
(Fisher OR~$= \foundKwFisherOR$, $p = \pFoundKwFisher$~\pFoundKwFisherSig{}).
The A++ sub-population (\NfoundKwTotalApp{} sites) shows a stronger
signal: \foundKwAppRate\% carry a founding keyword
(OR~$= \foundKwAppFisherOR$, $p = \pFoundKwAppFisher$~\pFoundKwAppFisherSig{}),
consistent with the most precisely aligned sites being disproportionately
historically foundational monuments.

The classifier assigns sites to five founding-related categories; full
keyword definitions and the site-by-site listing are in the supplementary
archive \citep{tenenbaum2026audit}.



%------------------------------------------------------------------------------

\begin{thebibliography}{99}
%------------------------------------------------------------------------------

\bibitem[Aveni(2001)]{aveni2001skywatchers}
Aveni, A.~F. (2001).
\textit{Skywatchers: A Revised and Updated Version of Skywatchers of
Ancient Mexico}.
University of Texas Press.

\bibitem[Aveni(2003)]{aveni2003archaeoastronomy}
Aveni, A.~F. (2003).
Archaeoastronomy in the ancient Americas.
\textit{Journal of Archaeological Research}, 11(2), 149--191.

\bibitem[Casson(1989)]{casson1989periplus}
Casson, L. (1989).
\textit{The Periplus Maris Erythraei: Text with Introduction, Translation,
and Commentary}.
Princeton University Press.

\bibitem[Ciolek(2004)]{ciolek2004owtrad}
Ciolek, T.~M. (2004).
\textit{Old World Trade Routes (OWTRAD) Project}.
Asia Pacific Research Online, Canberra.
\url{http://www.ciolek.com/OWTRAD/DATA/tmcOWTRAD.html}

\bibitem[Benjamini \& Hochberg(1995)]{benjamini1995controlling}
Benjamini, Y. \& Hochberg, Y. (1995).
Controlling the false discovery rate: a practical and powerful approach
to multiple testing.
\textit{Journal of the Royal Statistical Society, Series B}, 57(1), 289--300.

\bibitem[Crown(1989)]{crown1989}
Crown, A.~D. (1989).
\textit{The Samaritans}.
J.~C.~B. Mohr (Paul Siebeck).

\bibitem[Freeman(1976)]{freeman1976}
Freeman, P.~R. (1976).
A Bayesian analysis of the megalithic yard.
\textit{Journal of the Royal Statistical Society, Series A}, 139(1), 20--35.

\bibitem[Good(2005)]{good2005permutation}
Good, P.~I. (2005).
\textit{Permutation, Parametric, and Bootstrap Tests of Hypotheses}
(3rd ed.).
Springer.

\bibitem[Holm(1979)]{holm1979simple}
Holm, S. (1979).
A simple sequentially rejective multiple test procedure.
\textit{Scandinavian Journal of Statistics}, 6(2), 65--70.

\bibitem[Hansen(2012)]{hansen2012silkroad}
Hansen, V. (2012).
\textit{The Silk Road: A New History}.
Oxford University Press.

\bibitem[Harvey(1987)]{harvey1987mappamundi}
Harvey, P.~D.~A. (1987).
Medieval maps of the world.
In J.~B.~Harley \& D.~Woodward (Eds.),
\textit{The History of Cartography, Volume~1:
Cartography in Prehistoric, Ancient, and Medieval Europe
and the Mediterranean} (pp.~283--370).
University of Chicago Press.

\bibitem[Burton(1893)]{burton1893pilgrimage}
Burton, R.~F. (1893).
\textit{Personal Narrative of a Pilgrimage to Al-Madinah \& Meccah},
vol.~2 (Memorial Edition), p.~63, n.~2.
Tylston and Edwards.

\bibitem[Hunger \& Steele(2019)]{hunger2019mul}
Hunger, H. \& Steele, J. (2019).
\textit{The Babylonian Astronomical Compendium MUL.APIN}.
Routledge.

\bibitem[Thureau-Dangin(1921)]{thureau1921}
Thureau-Dangin, F. (1921).
\textit{Num\'eration et m\'etrologie sum\'eriennes}.
Revue d'Assyriologie et d'arch\'eologie orientale, 18, 123--142.

\bibitem[Kendall(1974)]{kendall1974}
Kendall, D.~G. (1974).
Hunting quanta.
\textit{Philosophical Transactions of the Royal Society A}, 276(1257), 231--266.

\bibitem[MacDonald(1964)]{macdonald1964}
MacDonald, J. (1964).
\textit{The Theology of the Samaritans}.
SCM Press.

\bibitem[Mardia \& Jupp(2000)]{mardia2000directional}
Mardia, K.~V. \& Jupp, P.~E. (2000).
\textit{Directional Statistics}.
Wiley.

\bibitem[Magli(2013)]{magli2013}
Magli, G. (2013).
\textit{Architecture, Astronomy and Sacred Landscape in Ancient Egypt}.
Cambridge University Press.

\bibitem[Magli(2016)]{magli2016}
Magli, G. (2016).
\textit{Archaeoastronomy: Introduction to the Science of Stars and Stones}.
Springer.

\bibitem[Miksic(1990)]{miksic1990borobudur}
Miksic, J.~N. (1990).
\textit{Borobudur: Golden Tales of the Buddhas}.
Periplus Editions.

\bibitem[Neugebauer(1975)]{neugebauer1975history}
Neugebauer, O. (1975).
\textit{A History of Ancient Mathematical Astronomy}, 3~vols.
Springer.

\bibitem[Powell(1987)]{powell1987metrological}
Powell, M.~A. (1987).
Masse und Gewichte.
In D. O. Edzard (Ed.),
\textit{Reallexikon der Assyriologie und Vorderasiatischen Arch\"aologie},
vol.~7, pp.~457--517. Walter de Gruyter.

\bibitem[Purvis(1968)]{purvis1968}
Purvis, J.~D. (1968).
\textit{The Samaritan Pentateuch and the Origin of the Samaritan Sect}.
Harvard University Press.

\bibitem[Robson(2008)]{robson2008mathematics}
Robson, E. (2008).
\textit{Mathematics in Ancient Iraq: A Social History}.
Princeton University Press.

\bibitem[Ruggles(1999)]{ruggles1999astronomy}
Ruggles, C.~L.~N. (1999).
\textit{Astronomy in Prehistoric Britain and Ireland}.
Yale University Press.

\bibitem[Ruggles \& Cotte(2010)]{ruggles2010icomos}
Ruggles, C.~L.~N. \& Cotte, M. (Eds.) (2010).
\textit{Heritage Sites of Astronomy and Archaeoastronomy in the Context
of the UNESCO World Heritage Convention: A Thematic Study}.
ICOMOS \& International Astronomical Union.

\bibitem[Ruggles(2015)]{ruggles2015handbook}
Ruggles, C.~L.~N. (Ed.) (2015).
\textit{Handbook of Archaeoastronomy and Ethnoastronomy}.
Springer.

\bibitem[Sen(2014)]{sen2014silk}
Sen, T. (2014).
The spread of Buddhism.
In B.~Z.~Kedar \& M.~E.~Wiesner-Hanks (Eds.),
\textit{The Cambridge World History, Vol.~5: Expanding Webs of Exchange
and Conflict, 500 CE--1500 CE} (pp.~447--479).
Cambridge University Press.

\bibitem[Silva(2020)]{silva2020statistical}
Silva, F. (2020).
A probabilistic framework and significance test for the analysis
of structural orientations in skyscape archaeology.
\textit{Journal of Archaeological Science}, 118, 105138.

\bibitem[Tenenbaum(2026)]{tenenbaum2026audit}
Tenenbaum, S. (2026).
\textit{Longitude Quantization in the UNESCO World Heritage Corpus: Domes, Stupas, and 3\textdegree{} Harmonics: Code and Data} (v1.6.0).
Zenodo. \url{https://doi.org/10.5281/zenodo.19938283}

\bibitem[Tomber(2008)]{tomber2008rome}
Tomber, R. (2008).
\textit{Indo-Roman Trade: From Pots to Pepper}.
Duckworth.

\bibitem[Thom(1964)]{thom1964}
Thom, A. (1964).
The larger units of length of megalithic man.
\textit{Journal of the Royal Statistical Society, Series~A}, 127(4), 527--533.

\bibitem[Thom(1967)]{thom1967}
Thom, A. (1967).
\textit{Megalithic Sites in Britain}.
Oxford University Press.

\bibitem[Thom(1971)]{thom1971}
Thom, A. (1971).
\textit{Megalithic Lunar Observatories}.
Oxford University Press.

\bibitem[UNESCO World Heritage Centre(2024)]{unescowhc2024xml}
{UNESCO World Heritage Centre}. (2024).
\textit{World Heritage List} [XML dataset].
Retrieved from \url{https://whc.unesco.org/en/list/xml/}.

\bibitem[Watson \& Horowitz(2011)]{watson2011writing}
Watson, R.~T. \& Horowitz, W. (2011).
\textit{Writing Science Before the Greeks: A Naturalistic Analysis of the
Babylonian Astronomical Treatise MUL.APIN}.
Brill.

\bibitem[Williamson \& Bellamy(1983)]{williamson1983leylines}
Williamson, T. \& Bellamy, L. (1983).
\textit{Ley Lines in Question}.
World's Work.

\end{thebibliography}


\section*{Data Availability Statement}
\label{sec:data_code}
All data and code are openly available; a permanent snapshot (v1.6.0) is
deposited at Zenodo (\citealt{tenenbaum2026audit};
\url{https://doi.org/10.5281/zenodo.19938283}); UNESCO XML: \citealt{unescowhc2024xml}.
Scripts, audit files, and maps are indexed in
Supplementary~S1 (\texttt{supplementary/Supplementary\_S1.md}), with each
file's role, macros produced, and paper section noted.
Full reproduction: \texttt{reproduce\_all\_macros.sh},
\texttt{update\_all\_audits.sh}, \texttt{generate\_figures.py}.

\section*{Acknowledgements}
The UNESCO World Heritage Centre maintains the publicly accessible XML
dataset on which this study depends. This research was conducted independently and without external institutional affiliation.

\section*{Conflict of Interest}
The author declares no conflict of interest.

\section*{Software and AI Disclosure}
All statistical analyses were performed in Python~3 using NumPy, SciPy, and pandas.
Large-language-model tools were used during preparation of this manuscript:
Claude (Anthropic) assisted with prose drafting and editing;
GitHub Copilot (Microsoft) assisted with code autocompletion.
These tools were used solely as writing and coding aids.
All scientific hypotheses, analytical choices, keyword-list decisions,
interpretation of results, and conclusions are solely the responsibility of the author.
No AI tool was used to generate, fabricate, or substitute for primary data or statistical analysis.
